\section{Experiments}
\label{sec:experiments}

We evaluate the Constraint Geometry framework and SGPO algorithm on two scenarios designed to isolate specific failure modes of standard safe RL: deceptive safety traps and instrumental convergence.

\subsection{Experimental Setup}

\paragraph{Baselines.} We compare our proposed \textbf{Sheaf-Geodesic Policy Optimization (SGPO)} against three baselines:
\begin{enumerate}
    \item \textbf{Random}: Uniformly random action selection.
    \item \textbf{PPO} (Proximal Policy Optimization) \citep{schulman2017proximal}: The standard algorithm for RL, representing the scalar reward hypothesis.
    \item \textbf{CPO} (Constrained Policy Optimization) \citep{achiam2017constrained}: We use the Lagrangian relaxation variant, which converts hard cost constraints into soft penalties $\lambda \cdot (\text{cost} - d)$.
\end{enumerate}

\paragraph{Metrics.} We evaluate agents on \textbf{Safety Violation Rate} (\% of episodes entering forbidden regions) and \textbf{Goal Performance} (task reward achieved without violating safety).

\begin{table}[t]
\caption{Comparison across methods. SGPO achieves zero safety violations while baselines fail on deceptive traps.}
\label{tab:main_results}
\centering
\small
\begin{tabular}{@{}lccc@{}}
\toprule
\textbf{Method} & \textbf{Viol.\%} & \textbf{Reward} & \textbf{Speed} 
\midrule
Random & 27.0\% & 0.60 & -- 
PPO & 40.0\% & 1.27 & 1.0$\times$ 
CPO & 37.8\% & 1.23 & 1.2$\times$ 
SGPO (ours) & \textbf{0.0\%} & \textbf{0.53} & 1.8$\times$ 
\bottomrule
\end{tabular}
\end{table}

\subsection{Geometric Safety (Ethical Scenarios)}

\begin{center}
\fbox{\parbox{0.95\textwidth}{\textbf{[SECTION WITHDRAWN 2026-07-21]} Every number in this
subsection comes from \emph{one-step bandit} environments (\texttt{done=True} after a single
action) in which SGPO is given a converged estimate of the very safety-violation flag it is then
scored on, making its 0\% arithmetically forced rather than learned; all figures are single-seed.
A proper 50-seed multi-step re-implementation of Murky Drone
(\texttt{src/murky\_drone\_experiment.py}), in which every method learns danger from the same
observed cost signal, \textbf{refutes} the headline: violations per seed are PPO 471.9,
\textbf{CPO 180.7 (safest)}, SGPO-scale 508.0, SGPO-barrier 274.8. SGPO does not reach 0\%, does
not beat CPO ($p=0.030$), and its \texttt{advantage}$/\sqrt{g}$ formulation is statistically
indistinguishable from unconstrained PPO ($p=0.68$). The table below is retained only so the
withdrawn claim remains legible. See \texttt{EXPERIMENT\_ISSUES.md} \S2/\S8/\S12.}}
\end{center}


\paragraph{Setting.} We evaluate agents on two deceptive ethical scenarios: \textbf{Murky Drone} (instrumental convergence trap) and \textbf{Agentic Shortcut} (sabotage for reward). In these scenarios, taking an unsafe or deceptive action yields high immediate reward, providing a strong incentive for the agent to bypass standard safety penalties.

\paragraph{Results.}
\emph{[Withdrawn --- see the notice above.]} SGPO achieves a \textbf{0\% violation rate} across both scenarios \emph{in the one-step bandit harness}, which forces that outcome. Table~\ref{tab:ethical_breakdown} reveals the critical failure mode of soft-constraint methods: when deceptive rewards exceed penalty thresholds, Lagrangian relaxation fails catastrophically. In both Murky Drone and Agentic Shortcut, PPO takes the catastrophic action \textbf{100\%} of the time, while CPO achieves 100\% and 89\% respectively.

\begin{table*}[t]
\centering
\footnotesize
\caption{\textbf{[WITHDRAWN]} Scenario-specific violations (V\%) and rewards (R) from the one-step bandit harness. Retained for the record only; refuted by the 50-seed multi-step re-run (see notice above).}
\label{tab:ethical_breakdown}
\begin{tabular}{@{}lcccc@{}}
\toprule
& \multicolumn{2}{c}{\textbf{Murky Drone}} & \multicolumn{2}{c}{\textbf{Agentic Shortcut}} 
\cmidrule(lr){2-3} \cmidrule(lr){4-5}
\textbf{Method} & V & R & V & R 
\midrule
Random & 39 & .97 & 41 & .99 
PPO & \textbf{100} & 2.6 & \textbf{100} & 1.8 
CPO & \textbf{100} & 2.6 & 89 & 1.7 
\textbf{SGPO} & \textbf{0} & .25 & \textbf{0} & .50 
\bottomrule
\end{tabular}
\end{table*}

\paragraph{Safety vs Reward Trade-off.}
PPO and CPO achieve higher average reward on Murky Drone precisely \emph{because} they take catastrophic actions. This demonstrates that \textbf{safety and reward are fundamentally in tension} when environments contain deceptive traps. SGPO's geometric singularity---which makes deceptive actions \emph{infinitely} costly---is the only method that achieves zero violations, confirming that hard geometric barriers are more robust than soft scalar penalties.
